\chapter{Bảng tra cứu, nguồn và tái lập}\label{app:reference}
\section{Các hằng số kernel}
Mọi kernel compact trong bảng dùng giá đỡ $[-1,1]$ và chuẩn hóa như
Chương 3. Với kernel Uniform có giá đỡ $[-1/2,1/2]$, các hằng số sẽ
khác; không trộn hai cách chuẩn hóa khi dùng cùng bandwidth.
\begin{center}
\begin{tabular}{lcc}
\toprule Kernel & $\mu_2(K)$ & $R(K)$\\\midrule
Gaussian & $1$ & $1/(2\sqrt\pi)$\\
Epanechnikov & $1/5$ & $3/5$\\
Uniform & $1/3$ & $1/2$\\
Triangular & $1/6$ & $2/3$\\
Biweight & $1/7$ & $5/7$\\
Triweight & $1/9$ & $350/429$\\
Tricube & $35/243$ & $175/247$\\\bottomrule
\end{tabular}
\end{center}
Các hằng số compact được kiểm chứng bằng tích phân đa thức chính xác
và tích phân số độc lập. Với Gaussian dùng tích phân chuẩn tắc.

\section{Bảng công thức và điều kiện tối thiểu cần nhớ}
\begin{longtable}{p{.20\textwidth}p{.72\textwidth}}
\toprule Chủ đề & Công thức và điều kiện\\\midrule\endhead
EDF & $\E\Fhat(x)=F(x)$, $\Var\Fhat(x)=F(x)(1-F(x))/n$;
mẫu i.i.d., không yêu cầu có mật độ.\\[4pt]
DKW & $\epsilon_n=\sqrt{\log(2/\alpha)/(2n)}$; băng cắt tại 0 và 1,
phủ đồng thời ít nhất $1-\alpha$.\\[4pt]
Plug-in & $T(\Fhat)$; nhất quán cần tính liên tục phù hợp hoặc LLN
và điều kiện moment của thống kê cụ thể.\\[4pt]
KS & $D_n=\max_i\{i/n-F_0(X_{(i)}),F_0(X_{(i)})-(i-1)/n\}$;
bảng null chuẩn dùng $F_0$ liên tục, đã chỉ định đầy đủ.\\[4pt]
Histogram & $N_j/(nh)$; bias thường $O(h)$, variance $O(1/(nh))$;
MISE $n^{-2/3}$ trong thiết lập trơn bậc một và phân hoạch thích hợp.\\[4pt]
KDE & $(nh)^{-1}\sum K((x-X_i)/h)$; $K\geq0$, $\int K=1$
đảm bảo mật độ.\\[4pt]
Bias KDE & $h^2\mu_2f''/2+O(h^4)$ dưới điều kiện đạo hàm bậc bốn
và moment tuyệt đối bậc bốn; không áp dụng ở biên chưa hiệu chỉnh.\\[4pt]
AMISE KDE & $h^4\mu_2^2R(f'')/4+R(K)/(nh)$;
cần điều kiện phần dư tích phân. $h_*\asymp n^{-1/5}$,
MISE $O(n^{-4/5})$.\\[4pt]
Normal reference & $(4/3)^{1/5}s n^{-1/5}$ với Gaussian và mật độ
tham chiếu chuẩn; scale bền vững hoặc hệ số $0.9$ cần ghi riêng.\\[4pt]
CV & $\int\widehat f_h^2-(2/n)\sum_i\widehat f_{(-i),h}(X_i)$;
giữ $h$, estimator LOO có mẫu số $n-1$.\\[4pt]
Phép chiếu & $\sum_{j\leq D}\widehat a_j\phi_j$;
rủi ro bằng sai số xấp xỉ cộng tổng variance hệ số. Không tự không âm.\\[4pt]
NW & $\sum k_iY_i/\sum k_i$; mẫu số dương, bias nội miền có
design bias $2r'f'/f$.\\[4pt]
Local linear & $(S_2T_0-S_1T_1)/(S_0S_2-S_1^2)$;
ma trận không suy biến, $\sum w_i=1$, $\sum w_i(X_i-x)=0$.\\[4pt]
Root--Unroot & $Y_j=\sqrt{(k/n)(N_j+1/4)}$;
variance gần $k/(4n)$ trong xấp xỉ Poisson đủ tốt;
hồi quy, lấy phần dương, bình phương và chuẩn hóa.\\[4pt]
Nhiều chiều & Với trơn bậc hai và xử lý biên phù hợp,
$h_*\asymp n^{-1/(4+d)}$, MISE $O(n^{-4/(4+d)})$.\\
\bottomrule
\end{longtable}

\section{Bản đồ nguồn để đọc sâu}
Số trang PDF tính từ 1 và có tính cả bìa; số trang in là số trong giáo
trình. Trong các slide, số slide trùng thứ tự trang PDF của các bản hiện có.
\begin{longtable}{p{.22\textwidth}p{.70\textwidth}}
\toprule Nội dung & Nguồn chính\\\midrule\endhead
EDF, plug-in, GOF & \cite{slide1}, 46 trang PDF; Wasserman Chương 2.
EDF/hội tụ đều ở slide 10--20; plug-in 30--35; GOF 36--46.\\[4pt]
Histogram & \cite{slide2}, slide 12--18; Wasserman Mục 6.2,
chứng minh MISE trong phụ lục trang in 140--142, PDF 149--151.\\[4pt]
KDE và bandwidth & \cite{slide2}, slide 19--29; Wasserman Mục 6.3,
trang in 132--136, PDF 141--145.\\[4pt]
Hölder và minimax & Tsybakov Mục 1.2, đặc biệt Mệnh đề 1.1--1.2
và Định lý 1.1, trang in 4--9, PDF 15--20;
Chương 2 cho công cụ lower bound.
Wasserman Định lý 6.31, trang in 134--135, PDF 143--144,
cho lower bound MISE trên lớp đạo hàm $L_2$.\\[4pt]
Phép chiếu & \cite{projection}, slide 8--24; \cite{slide2}, slide 30--34.
Tsybakov Mục 1.7 bổ sung nền chuỗi trực chuẩn; lưu ý phần đó chủ yếu
xét hồi quy, không thay nguyên xi variance hồi quy bằng variance mật độ.\\[4pt]
Hồi quy và băng & Wasserman Mục 5.4 về hồi quy cục bộ, Mục 5.7
về băng tin cậy; trang in 89--90, PDF 99--100, nêu vấn đề mục tiêu
làm trơn so với hàm thật. Tsybakov Mục 1.5--1.6.\\[4pt]
Root--Unroot & Wasserman Mục 6.6, trang in 139--140, PDF 148--149.\\[4pt]
KDE nhiều chiều & Wasserman Mục 6.5, trang in 138--139, PDF 147--148.\\
\bottomrule
\end{longtable}

\subsection{Các điểm được làm rõ khi đối chiếu}
\begin{enumerate}
\item Chệch histogram và KDE cần phân biệt giá trị tại điểm, toàn miền,
và tích phân. Không tích phân một chặn điểm hằng trên miền vô hạn.
\item Nhân $K_4=(3-u^2)\varphi/2$ là bậc bốn theo định nghĩa moment
của tài liệu; tên ``bậc ba'' trong một số quy ước chỉ đếm số moment
bị triệt tiêu. Công thức moment mới là căn cứ.
\item Hölder bậc lớn hơn 1 phải dùng đạo hàm; histogram không tự đạt
$D^{-2\beta}$ với mọi $\beta$ chỉ vì mật độ trơn.
\item Đồng nhất thức contrast dùng $-2\nu_n(g)$ với
$\nu_n=P_n-P$. Dấu được suy lại từ định nghĩa, không sao chép dấu cộng
trong phác thảo slide chiếu.
\item Chuyển mật độ từ $[0,1]$ về $[a,b]$ cần hệ số $1/(b-a)$;
$1/\sqrt{b-a}$ chỉ là hệ số chuẩn hóa cơ sở $L_2$.
\item Không gán một hằng số phạt phổ quát từ ví dụ mô phỏng cho một
định lý oracle. Phác thảo Chương 4 ghi tường minh bước tập trung được
viện dẫn.
\end{enumerate}

\section{Tái lập bản PDF và kiểm chứng}
Thư mục chứa \texttt{main.tex}, các chương trong \texttt{chapters/},
bảy hình PDF trong \texttt{figures/}, script \texttt{verify\_and\_figures.py}
và báo cáo \texttt{verification.json}.
Chạy tại thư mục tài liệu:
\begin{verbatim}
python3 verify_and_figures.py
latexmk -xelatex -interaction=nonstopmode -halt-on-error main.tex
\end{verbatim}
Python cần NumPy, SciPy, SymPy, Matplotlib. Trên môi trường workspace
hiện tại, interpreter có sẵn các thư viện là
\path{/opt/anaconda3/bin/python}; có thể thay \texttt{python3} bằng đường
dẫn đó. XeLaTeX dùng Latin Modern, \texttt{polyglossia} và các gói trong
preamble; không yêu cầu font thương mại hay font cài riêng từ hệ điều hành.
Lệnh \texttt{make check} chạy cả kiểm chứng số, biên dịch và kiểm tra cấu
trúc PDF bằng \texttt{check\_document.py}; bước cuối cần thêm PyMuPDF.
Nếu cần chọn interpreter, dùng
\texttt{make check PYTHON=/opt/anaconda3/bin/python}.

Script dùng seed cố định, kiểm tra moment bằng đại số và tích phân số,
kiểm tra công thức đạo hàm bên trong miền, giới hạn một phía ở bài gốc,
CV qua công thức và tích phân độc lập, KS qua thuật toán thư viện,
CvM/AD qua tích phân trực tiếp, và tái tạo tuyến tính ở cả hai biên.
Các phép kiểm số không thay thế chứng minh tiệm cận hoặc lower bound.
\intuition{Nguồn, quy ước và cách chạy là một phần của tính tái lập.
Nếu đáp số khác phần mềm, kiểm tra chuẩn hóa kernel, quantile, mẫu số
variance và quy ước đầu mút trước khi kết luận công thức sai.}
